\pagestyle{plain}

\chapter{CAPÍTULO 6: Design Participativo e Validação Territorial do Sistema Computacional de Salvaguarda de Saberes Agrícolas Tradicionais em Comunidades Quilombolas}


\begin{flushright}
Catuxe Varjão de Santana Oliveira  \footnote[1]{Programa de Pós-graduação em Ciências da Propriedade Intelectual-PPGPI/Universidade Federal de Sergipe-UFS. *E-mail:  \href{mailto:catuxe@academico.ufs.br}{\color{blue}{\underline{catuxe@academico.ufs.br}}}}
\\[0.3cm]
Luiz Diego Vidal Santos \footnote[2]{Universidade Estadual de Feira de Santana - UEFS, Brasil. E-mail: \href{mailto:ldvsantos@uefs.br}{\color{blue}{\underline{ldvsantos@uefs.br}}}}
\\[0.3cm]
Paulo Roberto Gagliardi \footnotemark[1]
\end{flushright}


% resumo em português — bookmark suprimido para manter hierarquia do capítulo
{\renewcommand{\PRIVATEbookmarkthis}[1]{}%
\begin{resumoumacoluna}
A eficácia operacional de sistemas computacionais de apoio à decisão em contextos comunitários está condicionada não apenas à robustez técnica de seus módulos, mas à incorporação de usuários finais na co-criação de interfaces, regras e fluxos de trabalho, assegurando aderência cultural, usabilidade e sustentabilidade institucional. Este estudo valida territorialemente o sistema computacional de salvaguarda de Saberes e Sistemas Agrícolas Tradicionais (SSAT) desenvolvido nos capítulos precedentes, mediante protocolo de Design Participativo (PD) conduzido junto a comunidades quilombolas do Território de Identidade Semiárido Nordeste II (Bahia). O processo participativo foi estruturado em cinco fases sequenciais. Na Fase 1 (Diagnóstico Participativo), grupos focais (n~=~3--4, 8--12 participantes) identificaram necessidades de documentação, desafios de salvaguarda e expectativas em relação a ferramentas computacionais. Na Fase 2 (Co-Criação), oficinas de co-design permitiram priorização comunitária de variáveis, calibração de limiares de funções de pertinência fuzzy, refinamento de regras de inferência e prototipagem de baixa fidelidade das interfaces. Na Fase 3 (Prototipagem Iterativa), um protótipo funcional em Python/Streamlit foi submetido a testes de usabilidade com tarefas dirigidas, mensuração de taxa de sucesso, tempo de conclusão e escala SUS, com ciclos iterativos até SUS $\geq$ 70. Na Fase 4 (Implantação Piloto), o sistema foi implantado em 2--3 comunidades com programa de capacitação (8--12h). Na Fase 5 (Avaliação de Adoção), métricas de uso, apropriação psicológica e impacto percebido foram avaliados após 6--12 meses.

\vspace{\onelineskip}
\noindent
\textbf{Palavras-chave}: Design participativo; Validação territorial; Usabilidade; Apropriação tecnológica; Sistema de apoio à decisão; Salvaguarda biocultural; Comunidades quilombolas.
\end{resumoumacoluna}}

\newpage

\section{Introdução}

O modelo integrado de decodificação e salvaguarda desenvolvido ao longo desta tese articulou quatro módulos sequenciais e interdependentes, produzindo um arcabouço capaz de converter saberes tácitos em variáveis auditáveis (Capítulo~3), sintetizar evidências sobre dimensões de vulnerabilidade (Capítulo~2), gerar o Índice de Salvaguarda Biocultural por integração TRI-fuzzy (Capítulo~4) e produzir dossiês com eficácia probatória validada por instâncias regulatórias (Capítulo~5). Contudo, a robustez técnica e a aceitação institucional não garantem, por si sós, a adoção efetiva e continuada do sistema pelas comunidades detentoras dos saberes, uma vez que a literatura sobre tecnologias para desenvolvimento comunitário documenta extensivamente a rejeição de ferramentas concebidas sem participação dos usuários finais \cite{Schuler1993,Simonsen2012}.

O Design Participativo (Participatory Design, PD) constitui uma abordagem centrada no usuário que promove o envolvimento ativo dos stakeholders ao longo de todo o ciclo de desenvolvimento de tecnologias, desde a definição de requisitos até a validação \cite{Schuler1993}. Para sistemas de apoio à decisão em contextos comunitários, o PD é imprescindível para assegurar aderência dos requisitos às necessidades reais, apropriação psicológica da tecnologia, capacitação para uso autônomo e sustentabilidade do sistema após o encerramento da intervenção de pesquisa \cite{Simonsen2012}. A teoria da apropriação tecnológica de \citeonline{Carroll2001} fornece o arcabouço analítico para compreender o processo pelo qual os usuários transitam de mero uso instrumental para sentimento de propriedade, adaptação criativa e evangelização da tecnologia.

A hipótese testada neste estudo sustenta que a incorporação dos usuários finais na etapa de design das regras de inferência do sistema computacional (participatory design) aumenta a apropriação psicológica da tecnologia de decodificação, resultando em taxa de adoção continuada superior à de ferramentas de gestão exógenas. Este capítulo responde à questão de pesquisa Q5, investigando como a incorporação de usuários finais na construção das regras e interfaces do sistema computacional afeta a apropriação, adoção e uso continuado da tecnologia de salvaguarda no território.


\section{Materiais e Métodos}

O protocolo de Design Participativo foi estruturado em cinco fases sequenciais, conforme detalhado na metodologia geral da tese.


\subsection{Fase 1: Diagnóstico Participativo de Necessidades}

Foram realizados grupos focais (n~=~3--4, com 8--12 participantes cada) nas comunidades quilombolas do Território de Identidade Semiárido Nordeste II (Bahia), com o objetivo de identificar necessidades de documentação, desafios na salvaguarda e proteção de saberes, e expectativas em relação a ferramentas computacionais de apoio. Os participantes foram selecionados por amostragem intencional, contemplando mestres de saberes, lideranças comunitárias, jovens aprendizes e representantes de associações quilombolas, assegurando diversidade de perspectivas geracionais e de gênero.

As técnicas empregadas compreenderam mapeamento de stakeholders, análise SWOT participativa, construção de personas (perfis típicos de usuários do sistema) e jornadas de usuário (\textit{user journeys}) descrevendo cenários de uso do sistema em contextos reais de documentação, proteção e tomada de decisão. As sessões foram gravadas em áudio/vídeo (com consentimento informado) e transcritas para análise temática \cite{Braun2006}.


\subsection{Fase 2: Co-Criação de Requisitos e Regras}

Oficinas de co-design (n~=~2--3, duração 4--6h cada) reuniram participantes das comunidades com a equipe de pesquisa para definição colaborativa de quatro elementos do sistema.

O primeiro elemento correspondeu à priorização de variáveis, onde os participantes avaliaram quais dimensões de vulnerabilidade e salvaguarda são mais relevantes para inclusão no ISB, complementando o painel Delphi (Capítulo~3) com a perspectiva comunitária direta. O segundo elemento envolveu a calibração de limiares, ajustando os intervalos de funções de pertinência fuzzy à semântica comunitária (por exemplo, definindo operacionalmente o que constitui ``alto risco de erosão intergeracional e migração'' na prática local). O terceiro elemento compreendeu o refinamento de regras, validando e ajustando regras de inferência do motor fuzzy mediante verificação de coerência com conhecimento local e identificação de regras ausentes ou contraditórias. O quarto elemento focou no design de interface, com prototipagem em papel (low-fidelity prototypes) das telas do sistema, definição de fluxo de navegação, terminologia acessível e visualizações de saída (gráficos radar de vulnerabilidade, rankings ISB, dossiês formatados).

As sessões empregaram cartões, diagramas impressos, sketches e post-its para facilitar ideação colaborativa, sendo gravadas em áudio/vídeo com consentimento para análise posterior.


\subsection{Fase 3: Prototipagem Iterativa e Testes de Usabilidade}

O protótipo funcional de média fidelidade foi implementado em Python com interface Streamlit, incorporando os requisitos co-criados na Fase 2. O protótipo integrou os módulos de entrada de dados (WOCAT-SLM-QBR), cálculo do ISB (módulo fuzzy), visualização de perfis de vulnerabilidade (gráficos radar) e geração de dossiês formatados para instâncias regulatórias (Capítulo~5).

O protocolo de teste de usabilidade envolveu 10--15 usuários executando tarefas dirigidas representativas dos cenários de uso identificados na Fase 1, incluindo geração do ISB para um saber específico, comparação de dois saberes por prioridade de salvaguarda e geração de dossiê para submissão ao IPHAN. As métricas quantitativas compreenderam taxa de sucesso nas tarefas, tempo de conclusão, número de erros e necessidade de assistência. As métricas qualitativas incluíram entrevistas pós-teste explorando dificuldades, sugestões de melhoria e satisfação geral mensurada pela System Usability Scale (SUS) \cite{Brooke1996}. Adicionalmente, análise \textit{think-aloud} foi conduzida, na qual os usuários verbalizaram pensamentos durante a execução de tarefas, permitindo identificar pontos de confusão e atrito.

Ciclos iterativos de redesign e teste foram conduzidos até que o protótipo alcançasse taxa de sucesso $\geq$ 85\% nas tarefas e escore SUS $\geq$ 70, correspondente ao limiar de usabilidade aceitável conforme classificação de \citeonline{Brooke1996}.


\subsection{Fase 4: Implantação Piloto e Capacitação}

O sistema validado na Fase 3 foi implantado em 2--3 comunidades quilombolas, acompanhado de programa de capacitação (oficinas de 8--12h totais) cobrindo três domínios. O primeiro domínio, conceitual, abordou fundamentos de salvaguarda de saberes tradicionais e lógica do sistema (sem formulações matemáticas formais). O segundo domínio, operacional, capacitou os participantes na entrada de dados, interpretação de resultados, geração de relatórios e dossiês. O terceiro domínio, estratégico, instruiu sobre integração do ISB em processos de proteção, geração de dossiês para instâncias regulatórias (IPHAN, INPI, CGen) e suporte a processos de registro e salvaguarda.

O suporte pós-implantação incluiu manual do usuário ilustrado, vídeos tutoriais curtos (3--5 minutos cada) e linha de suporte técnico via aplicativo de mensagens e telefone.


\subsection{Fase 5: Avaliação de Adoção e Apropriação}

Após período de 6--12 meses de uso autônomo, foi conduzida avaliação de adoção e apropriação mediante quatro dimensões.

\subsubsection{Métricas de uso}

Os logs do sistema registraram frequência de acesso, número de saberes documentados e classificados, funcionalidades mais utilizadas e padrões temporais de uso. A taxa de adoção foi definida como proporção de participantes da capacitação utilizando o sistema regularmente (ao menos uma vez por mês) após 6 meses, com meta de $\geq$ 60\%.

\subsubsection{Apropriação psicológica}

Foi aplicada a escala de apropriação tecnológica de \citeonline{Carroll2001}, adaptada ao contexto comunitário e validada linguisticamente, mensurar três dimensões: sentimento de propriedade (``o sistema é nosso''), adaptação criativa (uso de funcionalidades para fins não previstos) e evangelização (recomendação ativa a outras comunidades ou organizações). A meta estabelecida foi que $\geq$ 50\% dos usuários reportassem adaptações criativas ou evangelização ativa.

\subsubsection{Impacto percebido}

Entrevistas semiestruturadas e questionários exploraram se o sistema contribuiu para fortalecimento da documentação de saberes, proteção formal de saberes junto a instâncias regulatórias, argumentação em processos de licenciamento ou repartição de benefícios, e empoderamento na gestão de ativos bioculturais. A meta foi que $\geq$ 70\% dos usuários reportassem contribuição para pelo menos um resultado concreto.

\subsubsection{Sustentabilidade}

A capacidade de uso autônomo (sem suporte contínuo da equipe de pesquisa) foi avaliada pela proporção de tarefas completadas sem assistência externa, e a sustentabilidade institucional pela existência de rotinas comunitárias de alimentação e uso do sistema.


\subsection{Análise de Dados}

As métricas quantitativas de usabilidade (SUS, taxa de sucesso, tempo), adoção (frequência, taxa) e apropriação (escores da escala) foram analisadas por estatística descritiva (mediana, IQR), com comparações pré-pós conduzidas por teste de Wilcoxon para amostras pareadas quando aplicável. Os dados qualitativos (transcrições de grupos focais, entrevistas, sessões \textit{think-aloud}) foram analisados por codificação temática reflexiva \cite{Braun2006}, com triangulação entre fontes de dados para validação cruzada das categorias emergentes. A integração de dados foi conduzida por \textit{joint display} \cite{Creswell2018}, apresentando convergências e divergências entre achados quantitativos e qualitativos.


\section{Resultados}

\textit{[Os resultados serão apresentados após a execução do protocolo participativo. A estrutura prevista segue abaixo.]}

\subsection{Diagnóstico Participativo}

Os resultados dos grupos focais serão apresentados mediante mapa temático hierarquizado, organizando necessidades, desafios e expectativas em categorias e subcategorias emergentes. As personas e jornadas de usuário co-construídas serão apresentadas como artefatos do diagnóstico.

\subsection{Co-Criação e Refinamento do Sistema}

Os ajustes realizados no sistema a partir das oficinas de co-design serão detalhados, incluindo modificações nas funções de pertinência, regras de inferência adicionadas ou revisadas, e decisões de interface derivadas da prototipagem em papel. A convergência e divergência entre priorizações comunitárias e evidências meta-analíticas (Capítulo~2) serão discutidas.

\subsection{Usabilidade}

Os resultados dos testes de usabilidade serão reportados por ciclo iterativo, demonstrando a evolução dos escores SUS, taxas de sucesso e tempos de conclusão. Os achados da análise \textit{think-aloud} serão apresentados como categorias de atrito e suas resoluções ao longo dos ciclos.

\subsection{Adoção, Apropriação e Impacto}

As métricas de adoção (taxa de uso regular), apropriação (escores da escala de Carroll, evidências de adaptação criativa) e impacto percebido (resultados concretos reportados) serão apresentadas em formato integrado (\textit{joint display}).


\section{Discussão}

\textit{[A discussão será elaborada com base nos resultados obtidos. Os eixos analíticos previstos são apresentados abaixo.]}

Os resultados de adoção e apropriação serão interpretados à luz da teoria de apropriação tecnológica de \citeonline{Carroll2001} e comparados com taxas de adoção reportadas na literatura para tecnologias de gestão em contextos comunitários rurais. A relação entre intensidade da participação nas fases de co-criação e nível de apropriação será discutida como evidência do mecanismo causal proposto pela hipótese H5.

Os ajustes no motor fuzzy derivados da validação comunitária serão discutidos como contribuição para a robustez do ISB, demonstrando que a calibração participativa complementa (e não substitui) a calibração empírica e meta-analítica. As tensões entre critérios técnicos (TRI, meta-análise) e critérios comunitários serão analisadas como indicativas das fronteiras do design participativo em sistemas computacionais com forte componente técnico.

As limitações, incluindo o horizonte temporal de avaliação (6--12 meses), a cobertura geográfica restrita ao Semiárido Nordeste II, e o viés de seleção inerente à amostragem intencional nos grupos focais, serão explicitadas.


\section{Conclusão}

\textit{[A conclusão será elaborada após os resultados. As entregas previstas são apresentadas abaixo.]}

O protocolo de design participativo será apresentado como artefato replicável para validação territorial de sistemas computacionais de salvaguarda biocultural. Os indicadores de usabilidade (SUS), adoção (taxa de uso regular) e apropriação (sentimento de propriedade, adaptação criativa, evangelização) serão consolidados como evidência de que o sistema transcende o artefato técnico para constituir instrumento de empoderamento comunitário. A integração dos achados com os capítulos precedentes será sintetizada, demonstrando o fechamento do ciclo completo de decodificação, salvaguarda, auditoria e adoção territorial.


%% ==============================================
%% NOTAS PARA PREENCHIMENTO APÓS EXECUÇÃO
%% ==============================================
%% [GAP CRÍTICO: Após execução do protocolo participativo, inserir:
%%  - Mapa temático do diagnóstico participativo
%%  - Personas e user journeys co-construídas
%%  - Log de modificações no motor fuzzy por co-criação
%%  - Evolução dos escores SUS por ciclo iterativo
%%  - Taxas de sucesso e tempos de conclusão por tarefa
%%  - Métricas de adoção (frequência, taxa de uso regular)
%%  - Escores de apropriação (escala Carroll adaptada)
%%  - Joint display quantitativo-qualitativo
%%  - Evidências de adaptação criativa e evangelização]
